\section{Unidades utilizadas en las mediciones}
\label{sec:niveles}

\subsection{Decibeles absolutos y relativos}
\begin{itemize}
    \item \underline{dBm:} potencia relativa a $1 \, [mW]$, $P_{dBm} = 10 \log\left(P / 1 \, [mW]\right)$. Es una medida absoluta: $0 \, [dBm] = 1 \, [mW]$, $-30 \, [dBm] = 1 \, [\mu W]$ y $-60 \, [dBm] = 1 \, [nW]$.
    \item \underline{dB$\mu$V:} tensión relativa a $1 \, [\mu V]$. Como es un cociente de tensiones, se calcula como $V_{dB\mu V} = 20 \log\left(V / 1 \, [\mu V]\right)$. Es la unidad habitual en recepción de radiodifusión y en redes de cable, donde los niveles se especifican en tensión sobre $75 \, [\Omega]$.
    \item \underline{dBc:} nivel de una componente relativo al de la portadora, $N_{dBc} = N_{comp} - N_{portadora}$. Es una medida relativa: no depende de la impedancia ni de la ganancia del sistema, y por eso se usa para expresar bandas laterales, espurias, armónicos o ruido de fase.
\end{itemize}

\subsection{Relación entre potencia y tensión}
Sobre una resistencia $R$, una tensión eficaz $V$ disipa una potencia
\begin{equation}
    P = \frac{V^2}{R} \quad \Rightarrow \quad V = \sqrt{P R}
    \label{eq:p_v}
\end{equation}
Expresando $P$ en mW y $V$ en µV:
\begin{equation}
    V_{dB\mu V} = 20 \log\frac{\sqrt{P R}}{10^{-6}}
    = 10 \log\frac{P}{10^{-3}} + 10 \log\frac{10^{-3} R}{10^{-12}}
    = P_{dBm} + 90 + 10 \log R
    \label{eq:dbm_dbuv}
\end{equation}

Para los dos valores de impedancia habituales:
\begin{align}
    R = 50 \, [\Omega]: & \quad V_{dB\mu V} = P_{dBm} + 107{,}0 \\
    R = 75 \, [\Omega]: & \quad V_{dB\mu V} = P_{dBm} + 108{,}8
\end{align}
Por ejemplo, $-20 \, [dBm]$ equivalen a $87{,}0 \, [dB\mu V]$ ($22{,}36 \, [mV]$) sobre $50 \, [\Omega]$ y a $88{,}8 \, [dB\mu V]$ ($27{,}39 \, [mV]$) sobre $75 \, [\Omega]$. La diferencia entre ambos sistemas es $10 \log(75/50) = 1{,}76 \, [dB]$.

La \autoref{tab:conversion} resume algunas equivalencias.

\begin{table}[H]
    \centering
    \caption{Equivalencias entre dBm, potencia, tensión y dBµV para $50$ y $75 \, [\Omega]$.}
    \label{tab:conversion}
    \begin{tabular}{c c c c c c}
        \toprule
        \multirow{2}{*}{\textbf{dBm}} & \multirow{2}{*}{\textbf{Potencia}} & \multicolumn{2}{c}{\textbf{$50 \, [\Omega]$}} & \multicolumn{2}{c}{\textbf{$75 \, [\Omega]$}} \\
        \cmidrule(lr){3-4} \cmidrule(lr){5-6}
        & & \textbf{Tensión} & \textbf{dBµV} & \textbf{Tensión} & \textbf{dBµV} \\
        \midrule
        $0$   & $1 \, [mW]$   & $224 \, [mV]$  & $107{,}0$ & $274 \, [mV]$  & $108{,}8$ \\
        $-20$ & $10 \, [\mu W]$ & $22{,}4 \, [mV]$ & $87{,}0$ & $27{,}4 \, [mV]$ & $88{,}8$ \\
        $-40$ & $100 \, [nW]$ & $2{,}24 \, [mV]$ & $67{,}0$ & $2{,}74 \, [mV]$ & $68{,}8$ \\
        $-60$ & $1 \, [nW]$   & $224 \, [\mu V]$ & $47{,}0$ & $274 \, [\mu V]$ & $48{,}8$ \\
        $-107$ & $20 \, [fW]$ & $1{,}00 \, [\mu V]$ & $0{,}0$ & $1{,}22 \, [\mu V]$ & $1{,}8$ \\
        \bottomrule
    \end{tabular}
\end{table}

\subsection{Por qué hay que especificar la impedancia}
Un nivel en dBm indica una potencia y un nivel en dBµV, una tensión. Pasar de uno a otro requiere la \autoref{eq:p_v}, en la que interviene la resistencia sobre la que se desarrolla la tensión: la misma tensión de $1 \, [mV]$ entrega $20 \, [nW]$ ($-47 \, [dBm]$) sobre $50 \, [\Omega]$ y $13{,}3 \, [nW]$ ($-48{,}8 \, [dBm]$) sobre $75 \, [\Omega]$. Sin la impedancia, la conversión es ambigua en $1{,}76 \, [dB]$ entre los dos sistemas habituales.

En la práctica esto aparece al medir con el analizador de espectro, que tiene una entrada de $50 \, [\Omega]$ y mide en dBm, una red de $75 \, [\Omega]$ cuyos niveles se especifican en dBµV. Se pueden presentar tres casos:
\begin{itemize}
    \item \underline{Con adaptador de mínima pérdida:} a la lectura se le suma la atenuación del adaptador (unos $5{,}7 \, [dB]$) y se convierte con la ecuación de $75 \, [\Omega]$.
    \item \underline{Conexión directa:} se comete el error de desadaptación de $0{,}18 \, [dB]$ (\autoref{sec:equipamiento}).
    \item \underline{Analizador configurado en $75 \, [\Omega]$:} algunos instrumentos permiten indicar una impedancia de referencia de $75 \, [\Omega]$ y hacen la conversión internamente, pero eso no cambia la impedancia real de su entrada.
\end{itemize}

\section{Ejercicio de aplicación: índice de modulación de una señal AM}
\label{sec:ejercicio_am}

\textbf{Enunciado:} una señal AM ingresa al analizador de espectro con un nivel de portadora de $-20 \, [dBm]$ y bandas laterales de $-60 \, [dBm]$. Se requiere el índice de modulación $m$.

\textbf{Resolución en tensión:} se convierte cada nivel a tensión sobre $50 \, [\Omega]$ con la \autoref{eq:p_v}.
\begin{align}
    \text{Portadora:} \quad & P_c = 10^{-20/10} \, [mW] = 10 \, [\mu W]
    & E_c &= \sqrt{10^{-5} \cdot 50} = 22{,}36 \, [mV] \\
    \text{Banda lateral:} \quad & P_{BL} = 10^{-60/10} \, [mW] = 1 \, [nW]
    & E_{BL} &= \sqrt{10^{-9} \cdot 50} = 0{,}2236 \, [mV]
\end{align}

De la \autoref{eq:am}, cada banda lateral tiene amplitud $E_{BL} = \frac{m}{2} E_c$, de modo que
\begin{equation}
    m = \frac{2 E_{BL}}{E_c} = \frac{2 \cdot 0{,}2236}{22{,}36} = 0{,}02 = 2 \,\text{\%}
\end{equation}

\textbf{Resolución en dBc:} la banda lateral está $-60 - (-20) = -40 \, [dBc]$ respecto de la portadora. Con la \autoref{eq:indice_am}:
\begin{equation}
    m = 2 \cdot 10^{-40/20} = 2 \cdot 0{,}01 = 0{,}02 = 2 \,\text{\%}
\end{equation}

\textbf{Observaciones:}
\begin{itemize}
    \item El resultado no depende de la impedancia. Sobre $75 \, [\Omega]$ las tensiones son $27{,}39$ y $0{,}2739 \, [mV]$, pero su cociente es el mismo, porque la impedancia se cancela. Por eso el cálculo en dBc es el más directo.
    \item Un índice del $2 \,\text{\%}$ es muy bajo: las bandas laterales llevan solo $2 \cdot (m/2)^2 = 0{,}02 \,\text{\%}$ de la potencia de la portadora. Para que el método funcione, las bandas laterales tienen que quedar por encima del piso de ruido del analizador. Con $-60 \, [dBm]$ y un RBW de $9 \, [kHz]$ el margen es amplio, pero para índices menores conviene reducir el RBW.
    \item El cálculo supone modulación con un único tono sinusoidal, que es lo que permite identificar una única línea de banda lateral (\autoref{sec:am}).
\end{itemize}

\subsection{Ajustes recomendados del analizador}
Para las mediciones de la Parte 1, la cátedra recomienda:
\begin{itemize}
    \item Conector de desacople de continua, obligatorio.
    \item Medir solo señales de recepción.
    \item Nivel de referencia de $-30 \, [dBm]$, ajustable según la señal.
    \item \textit{Span} y frecuencias de inicio y fin según la banda medida.
    \item Ganancia de $5$ a $10 \, [dB]$.
    \item RBW de $9 \, [kHz]$ y VBW de $100 \, [kHz]$.
\end{itemize}
El RBW de $9 \, [kHz]$ es adecuado para explorar las bandas y para la FM y la televisión, pero es comparable al canal de AM de $10 \, [kHz]$, por lo que para resolver las bandas laterales de una emisora de AM conviene bajarlo a $1 \, [kHz]$ o menos (\autoref{sec:am}). El VBW, mayor que el RBW, deja la traza sin promediar. Si se lo reduce, la traza se suaviza y se lee mejor el nivel medio de una señal ruidosa como la OFDM.
